% !PRINTED 419
\weberpage{419}
%<<CONT
\weq{3}{\varPhi(x) = \varphi(x)\,\varphi(x^2)\,\varphi(x^3)\dots\varphi(x^{n-1})}

vanishes for $x = r$. But since $r$ may be any root of $X_n$, and
these roots are all different from one another, $\varPhi(x)$ must be
divisible by $X_n$, hence

\weq{4}{\varPhi(x) = X_n\,\varPsi(x).}

The function $\varPhi(x)$, as appears from its mode of formation, has the
coefficient $1$ for the highest power of $x$ and otherwise integer
coefficients; likewise $X_n$. Hence $\varPsi(x)$ must also have the
coefficient $1$ for the highest power of $x$ and otherwise integer
coefficients. But it follows from (2) and (3) that $\varPhi(x)$ takes
the value $\pm 1$ for $x = 1$, hence

\weq{5}{\pm 1 = X_n(1)\,\varPsi(1),}

wherein $\varPsi(1)$ must be an integer.

Now if $n$ is a prime power $p^{\pi}$, then $X_n(1) = p$, and this
contradicts equation (5). Thus the irreducibility of $X_n$ is proved
for the case in which $n$ is a prime or a power of a prime.

The general proof therefore still requires the following: \gesperrt{If $a$ and
$b$ are two integers without a common divisor, and if it is further assumed
that $X_a$ and $X_b$ are irreducible, it follows that $X_{ab}$
is also irreducible.}

To prove this, we set $n = ab$ and assume that $\varphi(x)$ is
a rational factor of $X_n$. The roots of $\varphi(x) = 0$ are then
contained among the primitive $n$th roots of unity, and are
therefore of the form

\weq{6}{\varrho = \alpha\beta,}

wherein $\alpha$ and $\beta$ denote certain primitive $a$th and
$b$th roots of unity (\S.\,132).

Now (by a special case of the Tschirnhausen transformation)
a function $\varPhi(x)$ with rational coefficients can be derived
from $\varphi(x)$, whose roots are the $a$th powers of the roots of
$\varphi(x)$. The roots of $\varPhi(x)$ are therefore
\[
\varrho^a = \beta^a,
\]
that is, since $a$ is relatively prime to $b$, all primitive $b$th
roots of unity; and if therefore any $\varrho$ exists at all,
$\varPhi(x)$ must have a non-constant factor in common with $X_b$. But now
we have assumed, as already proved, that $X_b$ has no rational
factor other than itself; hence $\varPhi(x)$ must be divisible by
%CONT>>
